\chapter{\ifenglish Methodology and System Design\else วิธีการดำเนินงานและการออกแบบระบบ\fi}

% =====================================================
\section{ภาพรวมของระบบ (System Overview)}

โครงงาน eDENT-ELDER Phase 2 ถูกพัฒนาขึ้นเพื่อแก้ไขข้อจำกัดของระบบคัดกรองรอยโรคในช่องปากแบบดั้งเดิม
โดยเฉพาะในบริบทของประเทศไทยที่มีข้อจำกัดด้านการเข้าถึงบริการทางการแพทย์
ในพื้นที่ชนบทและกลุ่มผู้สูงอายุ

จากการศึกษางานวิจัยในบทที่ 2 พบว่า
แม้ว่าจะมีการพัฒนา AI สำหรับวิเคราะห์ภาพทางการแพทย์อย่างต่อเนื่อง
แต่ระบบส่วนใหญ่ยังไม่สามารถนำไปใช้งานจริงได้ในระดับชุมชน
เนื่องจากขาดการออกแบบระบบที่รองรับ workflow ของผู้ใช้งานจริง
และไม่มีความเสถียรเพียงพอสำหรับการใช้งานในระดับ production

ดังนั้น โครงงานนี้จึงมุ่งเน้นการออกแบบระบบแบบ end-to-end
ที่สามารถผสาน Mobile Health, AI และ Telemedicine
เข้ากับการทำงานของบุคลากรสาธารณสุขได้อย่างมีประสิทธิภาพ

ระบบรองรับผู้ใช้งานหลายบทบาท ได้แก่
ผู้ป่วย อาสาสมัครสาธารณสุข (อสม.) ผู้คัดกรอง ผู้เชี่ยวชาญ และผู้ดูแลระบบ

📌 \textbf{[ใส่รูป]} System Overview Architecture

% =====================================================
\section{การออกแบบ Workflow ของระบบ (System Workflow)}

\subsection{เครื่องมือและเกณฑ์การคัดกรอง}

\begin{tcolorbox}[bluebox]
\textbf{Oral Health Assessment Tool (OHAT)}  
ใช้ประเมินความเสี่ยงของสภาพช่องปากเบื้องต้น

\textbf{Oral Health Impact Profile (OHIP)}  
ใช้ประเมินผลกระทบของสุขภาพช่องปากต่อคุณภาพชีวิต
\end{tcolorbox}

การเลือกใช้เครื่องมือมาตรฐานช่วยลดความแปรปรวนของผลลัพธ์
และเพิ่มความน่าเชื่อถือของระบบคัดกรอง

% -------------------------
\subsection{กระบวนการยืนยันตัวตน}

ระบบใช้การยืนยันตัวตนผ่าน OTP และ Consent ตาม PDPA
เพื่อให้มั่นใจว่าข้อมูลผู้ป่วยได้รับการจัดการอย่างปลอดภัย

\begin{enumerate}
\item กรอกหมายเลขโทรศัพท์
\item ยืนยัน OTP
\item ลง Consent
\item กำหนดบทบาท
\end{enumerate}

📌 \textbf{[ใส่รูป]} Authentication Flow

% -------------------------
\subsection{Workflow แบบ Multi-Role}

ระบบถูกออกแบบเป็น 5 เลน:

\begin{itemize}
\item Patient / VHV
\item AI System
\item Screener
\item Specialist
\item Administrator
\end{itemize}

📌 \textbf{[ใส่รูป]} 5-Lane Workflow Diagram

% -------------------------
\subsection{เหตุผลในการออกแบบ Workflow}

การออกแบบ workflow แบบ multi-role
ช่วยให้ระบบสอดคล้องกับการทำงานจริงในระบบสาธารณสุข

การใช้ asynchronous teleconsultation
ช่วยลดข้อจำกัดด้านเวลา และเพิ่ม scalability

% =====================================================
\section{สถาปัตยกรรมของระบบ (System Architecture)}

\subsection{ภาพรวมสถาปัตยกรรม}

📌 \textbf{[ใส่รูป]} High-Level Architecture

ระบบใช้แนวคิด Modular Monolith ร่วมกับ DDD
เพื่อแยก business logic ออกจาก infrastructure

% -------------------------
\subsection{Technology Stack}

\begin{table}[H]
\centering
\begin{tabular}{ll}
\toprule
Frontend & Next.js \\
Backend & FastAPI \\
Database & PostgreSQL \\
Storage & MinIO \\
Cache/Queue & Redis + Celery \\
\bottomrule
\end{tabular}
\caption{Technology Stack}
\end{table}

% -------------------------
\subsection{Low-Level Architecture}

ระบบ backend รองรับ:
\begin{itemize}
\item JWT Authentication
\item RBAC
\item Redis Queue
\item Audit Log (Hash-chain)
\end{itemize}

% -------------------------
\subsection{เหตุผลในการออกแบบ}

Modular Architecture ช่วย:
\begin{itemize}
\item ลด complexity
\item เพิ่ม maintainability
\item รองรับ scaling
\end{itemize}

% =====================================================
\section{การออกแบบระบบ AI (AI System Design)}

\subsection{Dual-Path AI}

📌 \textbf{[ใส่รูป]} Dual Path Diagram

\begin{itemize}
\item GPU (Primary)
\item CPU (Fallback)
\end{itemize}

% -------------------------
\subsection{เหตุผลในการออกแบบ}

เพื่อลด Single Point of Failure
และเพิ่ม High Availability

% =====================================================
\section{โครงสร้างพื้นฐานและ DevOps}

📌 \textbf{[ใส่รูป]} DevOps Pipeline

\begin{itemize}
\item Docker
\item CI/CD
\item Zero downtime deployment
\end{itemize}

% =====================================================
\section{ความปลอดภัยของระบบ}

\begin{itemize}
\item JWT
\item RBAC
\item Encryption
\item Audit Logging
\end{itemize}

ระบบสอดคล้องกับ PDPA และ HIPAA

% =====================================================
\section{สรุปวิธีการดำเนินงาน}

ระบบ eDENT-ELDER Phase 2
ถูกออกแบบให้ผสาน workflow จริงเข้ากับเทคโนโลยี

ส่งผลให้สามารถนำไปใช้งานจริงในระดับชุมชนได้
และมีความเสถียรในระดับ production